% Nội dung viết lại theo template đồ án game.
% Phạm vi yêu cầu hiện tại: chỉ mục 1.3 của Chương 1 và toàn bộ Chương 2.
% Dán các phần này vào đúng vị trí trong template Overleaf.

\section{Phạm vi và đối tượng người chơi}

Trò chơi trong đồ án được xây dựng theo hướng nhập vai phiêu lưu 2D kết hợp thu thập, huấn luyện và chiến đấu theo lượt. Phạm vi của sản phẩm tập trung vào trải nghiệm chơi đơn trên máy cá nhân, trong đó người chơi điều khiển một nhân vật chính di chuyển qua nhiều khu vực bản đồ, tương tác với nhân vật không điều khiển, nhận nhiệm vụ, gặp Pokemon hoang dã, tham gia các trận đấu, thu phục Pokemon, quản lý đội hình và tiếp tục tiến trình cốt truyện. Trò chơi không đặt mục tiêu trở thành sản phẩm thương mại hoàn chỉnh với số lượng nội dung lớn như các game nhập vai thương mại, mà tập trung chứng minh một vòng lặp gameplay đầy đủ và có khả năng mở rộng bằng dữ liệu cấu hình.

Về nội dung, phạm vi trò chơi bao gồm các khu vực overworld như nhà người chơi, thị trấn, phòng thí nghiệm, đường đi, hang động, thành phố theo hệ, gym và một số khu vực cao trào về cuối game. Người chơi có thể chuyển giữa các scene, nói chuyện với NPC, nhận vật phẩm từ rương, bước vào vùng cỏ để gặp Pokemon ngẫu nhiên, tương tác với Pokemon xuất hiện trực tiếp trên bản đồ và tham gia trainer battle. Hệ thống chiến đấu hỗ trợ hai dạng chính là wild battle và trainer battle. Wild battle cho phép người chơi bắt Pokemon bằng Pokeball, còn trainer battle tập trung vào việc đánh bại toàn bộ đội hình đối thủ để nhận thưởng, mở khóa tiến trình hoặc nhận badge.

Về hệ thống, phạm vi sản phẩm bao gồm những module cốt lõi của một game Pokemon-like RPG: điều khiển nhân vật, tương tác overworld, hội thoại, nhiệm vụ, cờ cốt truyện, battle theo lượt, dữ liệu Pokemon, dữ liệu chiêu thức, vật phẩm, shop, party, storage, Pokedex, save/load và main story sequence. Các module này đã được thiết kế để phối hợp với nhau thay vì tồn tại như các demo riêng lẻ. Ví dụ, một trận chiến có thể ảnh hưởng tới HP, kinh nghiệm, tiến hóa, trạng thái nhiệm vụ, Pokedex, vật phẩm, tiền thưởng, badge và dữ liệu lưu. Đây là phạm vi phù hợp với đồ án tốt nghiệp vì vừa đủ rộng để thể hiện năng lực thiết kế hệ thống, vừa đủ kiểm soát để có thể triển khai trong thời gian giới hạn.

Trò chơi hướng đến nhóm người chơi yêu thích game nhập vai 2D, đặc biệt là những người quen thuộc với cấu trúc khám phá bản đồ, đối thoại NPC, thu thập sinh vật và chiến đấu theo lượt. Nhóm người chơi chính có thể nằm trong độ tuổi từ khoảng 13 trở lên, vì game có cơ chế menu, chỉ số, hệ Pokemon, lựa chọn chiêu thức và quản lý đội hình đòi hỏi người chơi hiểu một số quy tắc cơ bản. Tuy nhiên, nội dung trò chơi không có yếu tố bạo lực đồ họa, máu me, ngôn ngữ nhạy cảm hoặc chủ đề người lớn, nên về định hướng phân loại nội dung có thể phù hợp với nhóm tuổi thiếu niên và người chơi phổ thông.

Nếu phân loại theo sở thích thể loại, trò chơi phù hợp nhất với người chơi thích RPG nhẹ, thích quá trình tăng trưởng nhân vật, thích sưu tầm và thích các trận đấu có tính chiến thuật vừa phải. Người chơi không cần phản xạ nhanh như trong game hành động thời gian thực, vì chiến đấu diễn ra theo lượt và mỗi lượt cho phép lựa chọn hành động. Nhờ vậy, trò chơi có thể tiếp cận cả người chơi casual muốn trải nghiệm câu chuyện và người chơi thích tối ưu đội hình Pokemon. Những người đã quen với các game Pokemon cổ điển sẽ nhanh chóng hiểu mục tiêu cơ bản của trò chơi, trong khi người chơi mới vẫn có thể làm quen thông qua các tình huống đơn giản như di chuyển, nói chuyện với NPC, bắt Pokemon đầu tiên và tham gia trận đấu đầu tiên.

Nếu phân loại theo nền tảng và hoàn cảnh chơi, trò chơi được thiết kế cho máy tính cá nhân, sử dụng bàn phím làm thiết bị điều khiển chính. Người chơi có thể chơi theo từng phiên ngắn hoặc dài nhờ hệ thống save/load. Điều này phù hợp với mô hình game nhập vai có tiến trình nhiều scene, vì người chơi không cần hoàn thành toàn bộ nội dung trong một lần chơi. Trò chơi cũng phù hợp để trình diễn trong môi trường học thuật, nơi người xem cần quan sát được các nhóm chức năng như di chuyển, tương tác, chiến đấu, nhiệm vụ, Pokedex và lưu dữ liệu trong một khoảng thời gian ngắn.

Về giới tính và nền văn hóa, trò chơi không định hướng riêng cho một nhóm giới tính cụ thể. Phong cách 2D, màu sắc tươi, nhân vật dạng hoạt hình, chiến đấu không mô tả bạo lực trực diện và chủ đề phiêu lưu thu thập giúp trò chơi có khả năng tiếp cận nhiều nhóm người chơi khác nhau. Về mặt văn hóa, trò chơi lấy cảm hứng từ cấu trúc Pokemon-like quen thuộc với cả người chơi Việt Nam và người chơi quốc tế. Tuy nhiên, vì phần hội thoại và thông báo trong game đang có nhiều nội dung tiếng Việt, phiên bản hiện tại phù hợp nhất với người chơi hiểu tiếng Việt. Nếu muốn mở rộng đối tượng, hệ thống cần bổ sung localization để chuyển đổi ngôn ngữ giao diện và lời thoại.

Phạm vi không bao gồm chế độ nhiều người chơi, trao đổi Pokemon qua mạng, chiến đấu online, hệ thống AI phức tạp như game thương mại, hoặc kho nội dung rất lớn. Những phần này đòi hỏi hạ tầng mạng, cân bằng dữ liệu và kiểm thử quy mô lớn, vượt quá mục tiêu hiện tại của đồ án. Thay vào đó, sản phẩm tập trung vào trải nghiệm chơi đơn có vòng lặp rõ ràng: khám phá, tương tác, chiến đấu, thu thập, phát triển đội hình, mở khóa cốt truyện và lưu tiến trình.

[CẦN BỔ SUNG] Nếu giảng viên yêu cầu phân loại giống ESRB/PEGI, có thể ghi: trò chơi dự kiến phù hợp nhóm tuổi 13+ do có yếu tố chiến đấu giả tưởng nhưng không có máu, không có nội dung tình dục, không có chất kích thích và không có ngôn ngữ thô tục. Nếu bạn muốn định vị mềm hơn, có thể đổi thành ``phù hợp người chơi phổ thông từ 10+ hoặc 12+'' tùy nội dung hội thoại thực tế.

\chapter{KHẢO SÁT VÀ PHÂN TÍCH YÊU CẦU}

Chương này trình bày khảo sát hiện trạng và phân tích yêu cầu của trò chơi theo định hướng của một đề tài game. Thay vì chỉ mô tả hệ thống dưới góc nhìn phần mềm, chương này tập trung làm rõ trò chơi được đặt trong bối cảnh nào, có điểm tương đồng và khác biệt gì so với các sản phẩm cùng thể loại, mục đích chơi là gì, người chơi sẽ sử dụng trò chơi trong tình huống nào và vì sao cách tiếp cận đó phù hợp. Các nội dung này là cơ sở để xây dựng phần thiết kế trò chơi, thiết kế gameplay, thiết kế kỹ thuật và đánh giá sản phẩm ở các chương sau.

\section{Khảo sát hiện trạng}

Thị trường và cộng đồng game đã có nhiều trò chơi nhập vai 2D dựa trên cấu trúc khám phá, nhiệm vụ, chiến đấu và phát triển nhân vật. Trong nhóm này, dòng game Pokemon là đại diện nổi bật cho mô hình thu thập sinh vật và chiến đấu theo lượt. Người chơi thường bắt đầu từ một khu vực an toàn, nhận Pokemon đầu tiên, di chuyển qua các tuyến đường, gặp Pokemon hoang dã, chiến đấu với trainer, đánh bại gym leader, nhận badge và mở khóa khu vực mới. Cấu trúc này dễ tiếp cận vì mỗi mục tiêu nhỏ đều rõ ràng, trong khi mục tiêu dài hạn là hoàn thiện đội hình và tiến xa hơn trong cốt truyện.

Các game Pokemon chính thức có ưu điểm ở hệ thống nội dung phong phú, số lượng Pokemon lớn, thế giới được thiết kế có tính nhận diện cao và luật chiến đấu đã được cân bằng qua nhiều thế hệ. Tuy nhiên, ở góc độ đồ án tốt nghiệp, việc tái tạo đầy đủ quy mô đó là không phù hợp. Một sản phẩm học thuật nên tập trung vào việc thiết kế và triển khai những cơ chế cốt lõi: hệ thống bản đồ, đối thoại, encounter, battle, item, party, quest, save/load và progression. Do đó, trò chơi trong đồ án không cạnh tranh bằng số lượng nội dung, mà thể hiện năng lực xây dựng nền tảng gameplay có thể mở rộng.

Ngoài Pokemon, một số game nhập vai 2D khác cũng có các đặc điểm liên quan. Các game JRPG cổ điển thường dùng bản đồ top-down, NPC, hội thoại, nhiệm vụ, trận đấu và tiến trình tuyến tính qua từng khu vực. Điểm mạnh của nhóm game này là khả năng kể chuyện và dẫn dắt người chơi qua thế giới được chia thành nhiều scene. Tuy nhiên, nhiều game JRPG tập trung vào đội hình nhân vật cố định, trong khi trò chơi của đồ án nhấn mạnh vào yếu tố thu thập Pokemon và thay đổi đội hình. Điều này tạo ra yêu cầu riêng về party, storage, Pokedex và dữ liệu từng cá thể Pokemon.

Một nhóm sản phẩm khác là các game monster-collecting độc lập, trong đó người chơi bắt và huấn luyện sinh vật để chiến đấu. Các game này thường phát triển thêm biến thể riêng về hệ, kỹ năng, tiến hóa hoặc cách bắt sinh vật. Điểm chung của chúng là vòng lặp thu thập và phát triển đội hình có sức hút dài hạn. Trò chơi trong đồ án kế thừa cách tiếp cận này ở mức cơ bản: mỗi Pokemon có base data, level, HP, move, hệ, tiến hóa và khả năng được bắt hoặc sở hữu. Điểm khác biệt là sản phẩm tập trung vào việc xây dựng hệ thống trong Unity bằng ScriptableObject và các manager runtime, phù hợp với mục tiêu kỹ thuật của đồ án.

Nếu so với các demo game 2D thường gặp trong học tập, sản phẩm có phạm vi rộng hơn. Một demo đơn giản có thể chỉ bao gồm di chuyển nhân vật, va chạm và một vài tương tác. Trong khi đó, trò chơi này kết nối nhiều module: PlayerController xử lý di chuyển, Overworld xử lý vùng cỏ và đối tượng bản đồ, BattleSystem xử lý trận đấu, QuestManager xử lý nhiệm vụ, MainStoryDirector xử lý step cốt truyện, Inventory xử lý vật phẩm, PokedexManager theo dõi Pokemon đã thấy/bắt, StorageSystem lưu Pokemon ngoài party và SaveLoadSystem lưu tiến trình. Việc nhiều module cùng tác động tới một vòng lặp chơi làm sản phẩm có độ phức tạp cao hơn một demo đơn lẻ.

Từ khảo sát trên, có thể rút ra một số yêu cầu quan trọng. Thứ nhất, game cần có vòng lặp chơi dễ hiểu để người chơi không bị lạc mục tiêu. Thứ hai, dữ liệu Pokemon, move, item và quest cần được tổ chức theo cách dễ mở rộng, vì số lượng nội dung có xu hướng tăng dần. Thứ ba, battle và overworld phải được liên kết chặt chẽ, vì kết quả chiến đấu ảnh hưởng tới nhiều trạng thái khác nhau. Thứ tư, save/load phải đủ đầy đủ để bảo toàn tiến trình của người chơi qua nhiều scene. Thứ năm, cốt truyện cần có cơ chế điều kiện như story flag để tránh sự kiện chạy sai thời điểm.

Sản phẩm trong đồ án có điểm mạnh là không chỉ mô phỏng một trận battle hoặc một bản đồ riêng lẻ, mà xây dựng một nền tảng gameplay tương đối hoàn chỉnh cho game Pokemon-like RPG. Các module đã có quan hệ thực tế với nhau: bắt Pokemon có thể cập nhật party, storage, Pokedex và quest; đánh trainer có thể trao thưởng, badge và mở cờ truyện; đi vào vùng cỏ có thể tạo battle; main story step có thể điều khiển hội thoại, NPC, battle và flag. Đây là điểm giúp sản phẩm có giá trị trình diễn và giá trị kỹ thuật trong bối cảnh đồ án.

[CẦN BỔ SUNG BẢNG] Có thể thêm bảng so sánh với 2--3 trò chơi hoặc nhóm trò chơi: Pokemon mainline, JRPG 2D cổ điển, monster-collecting indie game, demo Unity RPG. Các cột nên gồm: đặc điểm gameplay, cách xử lý battle, cách mở rộng nội dung, save/load, điểm liên quan tới đồ án.

\section{Mục đích của trò chơi}

Mục đích trực tiếp của trò chơi là mang lại trải nghiệm nhập vai 2D trong đó người chơi cảm nhận được quá trình trở thành một trainer: bắt đầu từ nhân vật có đội hình nhỏ, khám phá thế giới, gặp Pokemon, chiến đấu, thu phục, huấn luyện, nhận nhiệm vụ và vượt qua các thử thách theo tiến trình cốt truyện. Trò chơi hướng tới cảm giác tăng trưởng liên tục, khi mỗi trận đấu, mỗi Pokemon bắt được, mỗi item nhận được và mỗi badge đạt được đều làm người chơi thấy đội hình của mình mạnh hơn hoặc thế giới mở thêm lựa chọn mới.

Mục đích thứ hai là tạo một môi trường gameplay có tính khám phá. Người chơi không chỉ đi từ điểm bắt đầu tới điểm kết thúc, mà còn có thể kiểm tra NPC, rương, vùng cỏ, Pokemon ngoài map, shop và các khu vực mới. Việc chia thế giới thành nhiều scene như town, road, cave và gym giúp trò chơi có nhịp tiến triển rõ ràng. Mỗi khu vực có thể chứa encounter zone riêng, NPC riêng, trainer riêng và trigger cốt truyện riêng. Điều này làm người chơi có lý do để di chuyển và quan sát môi trường thay vì chỉ chọn trận đấu từ menu.

Mục đích thứ ba là tạo ra một hệ thống chiến đấu vừa dễ hiểu vừa có chiều sâu cơ bản. Battle theo lượt cho phép người chơi lựa chọn giữa đánh, đổi Pokemon, dùng item hoặc chạy khi được phép. Mỗi Pokemon có hệ và bộ chiêu thức, còn mỗi chiêu có power, accuracy, PP, category và hiệu ứng đặc biệt như gây trạng thái, tăng/giảm chỉ số hoặc hút máu. Người chơi không cần phản xạ nhanh, nhưng cần suy nghĩ về HP, move còn lại, loại Pokemon đối thủ và thời điểm dùng item. Đây là mức chiến thuật phù hợp với game nhập vai 2D học thuật.

Mục đích thứ tư là phục vụ việc chứng minh năng lực thiết kế hệ thống game. Với tư cách là đồ án tốt nghiệp, trò chơi không chỉ được đánh giá như một sản phẩm giải trí mà còn là một sản phẩm phần mềm. Vì vậy, mục đích kỹ thuật của trò chơi là thể hiện cách tổ chức dữ liệu, cách điều phối state, cách chuyển giữa overworld và battle, cách lưu tiến trình, cách viết hệ thống quest theo event và cách dùng story flag để kiểm soát cốt truyện. Những mục đích này giúp sản phẩm có giá trị học thuật, không chỉ là một bản chơi thử.

Trò chơi không đặt mục tiêu giáo dục một kiến thức bên ngoài như toán học, sinh học hay ngoại ngữ. Tuy nhiên, nó vẫn có giá trị rèn luyện tư duy hệ thống và ra quyết định trong phạm vi gameplay. Người chơi cần hiểu quy tắc, đọc phản hồi, quản lý tài nguyên, chọn hành động phù hợp và theo dõi tiến trình. Nếu dùng trong môi trường học tập lập trình game, sản phẩm có thể đóng vai trò ví dụ cho cách xây dựng một game RPG 2D có nhiều module liên kết với nhau.

\section{Định hướng sử dụng}

Trò chơi được định hướng sử dụng như một sản phẩm chơi đơn độc lập trên máy tính cá nhân. Người chơi có thể khởi động game, bắt đầu hoặc tiếp tục tiến trình, di chuyển trong thế giới game, tham gia các tình huống cốt truyện và lưu lại trạng thái khi cần dừng. Cách sử dụng này phù hợp với bản chất của game nhập vai, vì người chơi cần thời gian để tích lũy Pokemon, level, item, badge và tiến độ nhiệm vụ qua nhiều phiên chơi.

Trong bối cảnh đồ án, trò chơi cũng được định hướng sử dụng như một sản phẩm trình diễn kỹ thuật. Người đánh giá có thể quan sát một số luồng tiêu biểu: bắt đầu game, xem prologue, nhận Pokemon đầu tiên, đi vào vùng cỏ, vào battle, dùng item, bắt Pokemon, nói chuyện với NPC, nhận quest, đánh trainer, nhận badge, mở Pokedex và lưu/tải game. Các luồng này cho thấy hệ thống không chỉ có giao diện mà có trạng thái nội bộ và dữ liệu được cập nhật xuyên suốt.

Nếu sử dụng trong một buổi demo, người trình bày nên chuẩn bị save file hoặc scene bắt đầu phù hợp để rút ngắn thời gian. Ví dụ, một save ở giai đoạn đầu có thể dùng để trình diễn hệ thống starter, encounter và battle cơ bản; một save ở giữa game có thể dùng để trình diễn gym, badge, Pokemon đã tiến hóa và story flag; một save gần cuối có thể dùng để trình diễn trainer mạnh, nhiều scene và Pokedex. Cách chuẩn bị này giúp người xem thấy được nhiều nhóm chức năng mà không cần chơi toàn bộ từ đầu.

Trò chơi cũng có thể được dùng làm nền tảng phát triển tiếp. Vì dữ liệu Pokemon, move, item và quest được tổ chức bằng asset cấu hình, người phát triển có thể thêm nội dung mới mà không cần thay đổi quá nhiều logic lõi. Ví dụ, để thêm một vùng cỏ mới, người phát triển có thể tạo EncounterZone và danh sách WildPokemon; để thêm một đoạn cốt truyện, có thể thêm MainStoryStep và các action tương ứng; để thêm vật phẩm, có thể tạo ItemBase và đưa vào inventory/shop. Nhờ vậy, định hướng sử dụng không chỉ là chơi thử mà còn là tiếp tục mở rộng nội dung.

Trong phiên bản hiện tại, trò chơi chưa định hướng sử dụng trong môi trường nhiều người chơi hoặc triển khai online. Người chơi không cần tài khoản, không cần kết nối mạng và không có trao đổi dữ liệu với server. Đây là lựa chọn phù hợp với phạm vi đồ án vì giảm rủi ro liên quan tới đồng bộ mạng, bảo mật, cân bằng PvP và lưu trữ cloud. Toàn bộ trọng tâm được đặt vào vòng lặp gameplay local và cấu trúc hệ thống trong Unity.

\section{Cơ sở của định hướng sử dụng}

Cơ sở đầu tiên của định hướng sử dụng độc lập là tính phù hợp giữa game nhập vai 2D và mô hình chơi theo phiên. Game nhập vai thường có tiến trình dài hơn một màn chơi ngắn, nên cần save/load để người chơi có thể tiếp tục sau khi thoát. Project hiện đã có SaveLoadSystem lưu nhiều loại dữ liệu như party, storage, tiền, scene, vị trí người chơi, quest snapshot, story flag, Pokedex, inventory, NPC state, trigger đã chạy và Pokemon ngoài map đã bắt. Điều này hỗ trợ trực tiếp cho định hướng chơi đơn nhiều phiên.

Cơ sở thứ hai là vòng lặp Pokemon-like có khả năng duy trì động lực người chơi. Vòng lặp khám phá, gặp Pokemon, battle, bắt Pokemon, nhận kinh nghiệm, mở nhiệm vụ và tiến sang khu vực mới tạo ra nhiều tầng mục tiêu. Người chơi có mục tiêu ngắn hạn như thắng một trận hoặc bắt một Pokemon, mục tiêu trung hạn như xây dựng đội hình và mục tiêu dài hạn như hoàn thành cốt truyện hoặc đánh bại các gym. Sự kết hợp này giúp trò chơi tránh cảm giác chỉ lặp lại một hành động đơn giản.

Cơ sở thứ ba là tính phù hợp của battle theo lượt với đối tượng người chơi phổ thông. So với battle hành động thời gian thực, battle theo lượt giảm yêu cầu về phản xạ và tăng vai trò của lựa chọn chiến thuật. Người chơi có thời gian quan sát HP, move, item và trạng thái trước khi quyết định. Cách chơi này phù hợp với cả người chơi casual lẫn người chơi thích tối ưu đội hình. Nó cũng phù hợp với môi trường demo vì người xem có thể dễ dàng theo dõi từng hành động và kết quả.

Cơ sở thứ tư là khả năng mở rộng nội dung bằng dữ liệu cấu hình. Một trò chơi nhập vai sẽ khó phát triển nếu mọi Pokemon, move, item và nhiệm vụ đều bị viết cứng trong code. Project hiện sử dụng ScriptableObject cho PokemonBase, MoveBase, ItemBase và Quest, đồng thời sử dụng MainStorySequence để mô tả các bước cốt truyện. Cách tổ chức này tạo điều kiện để trò chơi được mở rộng theo hướng thêm nội dung thay vì viết lại logic. Đây là nền tảng hợp lý cho một sản phẩm đồ án vì có thể chứng minh thiết kế lâu dài, không chỉ chứng minh một luồng chơi cố định.

Cơ sở thứ năm là sự phù hợp giữa định hướng trình diễn kỹ thuật và cấu trúc module hiện có. Project được chia thành nhiều module như Battle, Pokemon, Player, Overworld, GamePlay, Menu và Core. Mỗi module phụ trách một nhóm vấn đề cụ thể. Khi trình diễn, người xem có thể liên hệ trực tiếp chức năng trên màn hình với phần thiết kế bên dưới: battle tương ứng với BattleSystem, dữ liệu Pokemon tương ứng với PokemonBase/Pokemon, nhiệm vụ tương ứng với QuestManager, cốt truyện tương ứng với MainStoryDirector, lưu dữ liệu tương ứng với SaveLoadSystem. Điều này giúp báo cáo và sản phẩm có tính nhất quán.

Cơ sở cuối cùng là mức độ nội dung hiện tại đủ để thể hiện định hướng trò chơi. Project có nhiều scene, nhiều Pokemon asset, nhiều move asset, hệ thống item, shop, party, storage, Pokedex, story flags và main story asset. Dù chưa đạt mức hoàn thiện của game thương mại, số lượng thành phần này đủ để chứng minh rằng định hướng sử dụng là một game nhập vai 2D có tiến trình, không phải một prototype chỉ có một tính năng. Trong giai đoạn tiếp theo, việc bổ sung hình ảnh minh họa, bảng so sánh, biểu đồ game flow và kết quả playtest sẽ giúp phần khảo sát và phân tích yêu cầu thuyết phục hơn.

Kết chương, trò chơi được xác định là một game nhập vai 2D chơi đơn, hướng tới người chơi yêu thích khám phá, thu thập Pokemon và chiến đấu theo lượt. Khảo sát hiện trạng cho thấy sản phẩm kế thừa cấu trúc quen thuộc của các game Pokemon-like và JRPG 2D, nhưng tập trung vào phạm vi học thuật là xây dựng nền tảng gameplay có nhiều module liên kết. Mục đích, định hướng sử dụng và cơ sở của định hướng này sẽ là căn cứ cho phần thiết kế trò chơi ở các chương tiếp theo.
